\documentclass[11pt]{article}
\usepackage{mathunal-base}
\mathunalfuente{serif}
\newcommand{\rp}{\underline{\hspace{2.2cm}}}

\newcommand{\temacuerpo}[8]{%
% #1 coef Q1  #2 rhs Q1  #3 y0 Q1  #4 matriz Q2  #5 Phi Q3a  #6 b) enunciado  #7 c) corte  #8 d) F(s)
\begin{enumerate}
\item $(30\%)$ Halle la solución del problema de valor inicial
\[
\begin{cases}
y'(t) + #1\, y(t) - #1 \displaystyle\int_0^t y(\tau)\operatorname{sen}(t - \tau)\,d\tau = #2\operatorname{sen}(t), \\[4pt]
y(0) = #3.
\end{cases}
\]
\textbf{Nota:} debe justificar paso a paso su procedimiento.

\item $(30\%)$ La solución general del sistema homogéneo de ecuaciones diferenciales lineales de primer orden
\[
\mathbf{X}'(t) = #4 \mathbf{X}(t)
\]
en el intervalo $(-\infty, +\infty)$ es
\[
\mathbf{X}(t) = c_1 \begin{pmatrix} a \\ 1 \end{pmatrix} e^{rt}
+ c_2 \left[ \begin{pmatrix} a \\ 1 \end{pmatrix} t e^{rt} + \begin{pmatrix} c \\ 0 \end{pmatrix} e^{rt} \right],
\]
donde $\lambda = \rp$, \quad $a = \rp$, \quad $c = \rp$.
\textbf{Nota:} debe justificar paso a paso su procedimiento.

\item $(40\%)$ En los literales a), b), c) y d) complete según corresponda. Sólo se calificará la respuesta; no es necesario que escriba el procedimiento.
\begin{enumerate}[label=(\alph*)]
\item $(10\%)$ Considere el sistema homogéneo $\mathbf{X}'(t) = \mathbf{A}\mathbf{X}(t)$, donde $\mathbf{A} \in \mathbb{R}_{2\times 2}$. Se sabe que
\[
\Phi(t) = #5
\]
es una matriz fundamental del sistema en $(-\infty, \infty)$. Los valores propios de $\mathbf{A}$ son $\lambda_1 = \rp$ y $\lambda_2 = \rp$. Un vector propio de $\mathbf{A}$ asociado a $\lambda_2$ es $\mathbf{K} = \rp$.
\item $(10\%)$ #6 \\
donde $a = \rp$, \ $b = \rp$, \ $c = \rp$, \ $f = \rp$, \ $g = \rp$, \ $h = \rp$.
\item $(10\%)$ Sea $f(t) = \begin{cases} 0 & \text{si } 0 \leq t < #7, \\ e^t & \text{si } #7 \leq t, \end{cases}$ \ entonces $\mathcal{L}\{f(t)\} = \rp$.
\item $(10\%)$ Si $F(s) = \ln\!\left(#8\right)$, entonces su transformada inversa es $\mathcal{L}^{-1}\{F(s)\} = \rp$.
\end{enumerate}
\end{enumerate}
}

\begin{document}

{\large\bfseries Tercer Parcial de Ecuaciones Diferenciales --- 28 de noviembre de 2022}

\medskip
\noindent Escuela de Matemáticas. Universidad Nacional de Colombia, Sede Medellín.

\medskip
\noindent\textit{Duración: 1 hora y 50 minutos. Tres preguntas. Puntaje: 1 (/30), 2 (/30), 3 (/40). En los puntos 1 y 2 debe justificar paso a paso su procedimiento; en el punto 3 sólo se califica la respuesta. No está permitido hacer preguntas, usar calculadoras, sacar hojas en blanco ni ningún tipo de apuntes. Este documento reúne los cuatro temas (A, B, C y D).}

\bigskip\hrule\medskip
\begin{center}\bfseries\large Tema A\end{center}
\temacuerpo{1}{-}{1}%
{\begin{pmatrix} 4 & -7 \\[2pt] \tfrac{7}{4} & -3 \end{pmatrix}}%
{\begin{pmatrix} e^{3t} & e^{7t} \\[2pt] -3e^{3t} & 9e^{7t} \end{pmatrix}}%
{Se sabe que $\mathbf{A} \in \mathbb{R}_{2\times 2}$ es una matriz de entradas constantes que satisface
$\mathbf{A}\begin{pmatrix} -2+3i \\ 3+2i \end{pmatrix} = (1+3i)\begin{pmatrix} -2+3i \\ 3+2i \end{pmatrix}$.
La solución general del sistema $\mathbf{X}'(t) = \mathbf{A}\mathbf{X}(t)$ en $(-\infty, \infty)$ es
$\mathbf{X}(t) = c_1 e^{bt}\begin{pmatrix} a\cos(ct) - 3\sin(ct) \\ 3\cos(ct) + f\sin(ct) \end{pmatrix}
+ c_2 e^{bt}\begin{pmatrix} g\cos(ct) + h\sin(ct) \\ 2\cos(ct) + 3\sin(ct) \end{pmatrix}$,}%
{3}{\dfrac{3s+7}{-5+2s}}

\bigskip\hrule\medskip
\begin{center}\bfseries\large Tema B\end{center}
\temacuerpo{\tfrac{1}{3}}{-2}{2}%
{\begin{pmatrix} -2 & -\tfrac{7}{4} \\[2pt] 7 & 5 \end{pmatrix}}%
{\begin{pmatrix} -2e^{-t} & 5e^{2t} \\[2pt] 3e^{-t} & -4e^{2t} \end{pmatrix}}%
{Se sabe que $\mathbf{A} \in \mathbb{R}_{2\times 2}$ es una matriz de entradas constantes que satisface
$\mathbf{A}\begin{pmatrix} 3-5i \\ 5+3i \end{pmatrix} = (2+4i)\begin{pmatrix} 3-5i \\ 5+3i \end{pmatrix}$.
La solución general del sistema $\mathbf{X}'(t) = \mathbf{A}\mathbf{X}(t)$ en $(-\infty, \infty)$ es
$\mathbf{X}(t) = c_1 e^{bt}\begin{pmatrix} a\cos(ct) + 5\sin(ct) \\ 5\cos(ct) + f\sin(ct) \end{pmatrix}
+ c_2 e^{bt}\begin{pmatrix} g\cos(ct) + h\sin(ct) \\ 3\cos(ct) + 5\sin(ct) \end{pmatrix}$,}%
{5}{\dfrac{3s-6}{5+2s}}

\bigskip\hrule\medskip
\begin{center}\bfseries\large Tema C\end{center}
\temacuerpo{\tfrac{1}{4}}{-3}{3}%
{\begin{pmatrix} 3 & 4 \\[2pt] -1 & 7 \end{pmatrix}}%
{\begin{pmatrix} e^{-3t} & e^{2t} \\[2pt] -4e^{-3t} & e^{2t} \end{pmatrix}}%
{Se sabe que $\mathbf{A} \in \mathbb{R}_{2\times 2}$ es una matriz de entradas constantes que satisface
$\mathbf{A}\begin{pmatrix} 2 \\ 1+2i \end{pmatrix} = (-2-3i)\begin{pmatrix} 2 \\ 1+2i \end{pmatrix}$.
La solución general del sistema $\mathbf{X}'(t) = \mathbf{A}\mathbf{X}(t)$ en $(-\infty, \infty)$ es
$\mathbf{X}(t) = c_1 e^{bt}\begin{pmatrix} a\cos(ct) \\ \cos(ct) + f\sin(ct) \end{pmatrix}
+ c_2 e^{bt}\begin{pmatrix} g\cos(ct) + h\sin(ct) \\ 2\cos(ct) + \sin(ct) \end{pmatrix}$,}%
{7}{\dfrac{2s+1}{3+2s}}

\bigskip\hrule\medskip
\begin{center}\bfseries\large Tema D\end{center}
\temacuerpo{\tfrac{1}{5}}{-4}{4}%
{\begin{pmatrix} -1 & -\tfrac{3}{2} \\[2pt] 6 & 5 \end{pmatrix}}%
{\begin{pmatrix} \sqrt3\,e^{2t} & e^{-2t} \\[2pt] e^{2t} & -\sqrt3\,e^{-2t} \end{pmatrix}}%
{Se sabe que $\mathbf{A} \in \mathbb{R}_{2\times 2}$ es una matriz de entradas constantes que satisface
$\mathbf{A}\begin{pmatrix} 7-2i \\ 3+5i \end{pmatrix} = (-4-5i)\begin{pmatrix} 7-2i \\ 3+5i \end{pmatrix}$.
La solución general del sistema $\mathbf{X}'(t) = \mathbf{A}\mathbf{X}(t)$ en $(-\infty, \infty)$ es
$\mathbf{X}(t) = c_1 e^{bt}\begin{pmatrix} a\cos(ct) + 2\sin(ct) \\ 3\cos(ct) + f\sin(ct) \end{pmatrix}
+ c_2 e^{bt}\begin{pmatrix} g\cos(ct) + h\sin(ct) \\ 7\cos(ct) + 3\sin(ct) \end{pmatrix}$,}%
{9}{\dfrac{2s-4}{3-s}}

\vfill
\noindent{\small\itshape Transcripción digital --- MathUNAL}

\end{document}
